\documentclass[10pt,a4paper]{amsart}
\usepackage[margin=19mm]{geometry}
\usepackage{amsmath,amssymb,mathtools}
\usepackage[scr=rsfs,cal=euler]{mathalfa}
\usepackage{xfrac}
\usepackage[unicode,hidelinks,bookmarks=false]{hyperref}
\newcommand{\ie}{i.e.\ }
\newcommand{\cf}{cf.\ }
\newcommand{\resp}{respectively\ }
\newcommand{\Subsection}{\subsection}
\providecommand{\nobreakdash}{\nobreak\hbox{-}}
\setlength{\emergencystretch}{2em}
\multlinegap0pt
\usepackage{bm}
\usepackage{bbm}
\usepackage{dsfont}
\usepackage{xspace}
\usepackage{cancel}
\usepackage{mathtools}
\usepackage[shortlabels]{enumitem}
\usepackage{amsthm}
\makeatletter
\@ifundefined{corollary}{\newtheorem{corollary}{\hskip\parindent Corollary}}{}
\makeatother
\renewcommand\Im{\operatorname{Im}}
\renewcommand\Re{\operatorname{Re}}
\title[Classification of the real Painlev\'{e} I transcendents by z]{Classification of the real Painlev\'{e} I transcendents by zeros and connection problem: an asymptotic study}
\author{Yan Huang, Yu-Tian Li, Wen-Gao Long}
\date{Source: arXiv:2506.03648v1 (CC BY 4.0)}
\begin{document}
\maketitle

\noindent\emph{Source and licence.} Classification of the real Painlev\'{e} I transcendents by zeros and connection problem: an asymptotic study. Version arXiv:2506.03648v1 (CC BY 4.0), distributed under the Creative Commons Attribution 4.0 International licence, as recorded in the arXiv licence metadata for this exact version and shown on the abstract page https://arxiv.org/abs/2506.03648v1. Full rights notice: licenses/arxiv-2506.03648-CC-BY-4.0.txt.\par\smallskip

\noindent\textit{Editorial scope.} Counted unit: the Corollary whose source label is \texttt{cor-classification-1} in the article (source0.tex lines 483--496, source label \texttt{cor-classification-1}), delivered together with the article's own complete original proof of it (source0.tex lines 499--513, one \texttt{proof} environment of 15 lines). No mathematics is written, repaired or re-derived: statement and proof are the article's own text line for line. Required context retained uncounted: eq-classifed-by-stokes (lines 142--148); eq-stokes-full-asymptotic-case-II (lines 420--427). Every definition, display and label of those lines is retained unchanged. Omitted from this package and not redistributed: 1--141, 149--419, 428--482, 497--498, 514--1295. Nothing inside a retained excerpt is omitted or altered: every retained body is delivered exactly as the article's lines are written, so no omission receipt of any kind accompanies this package. Citations: the retained excerpts carry no citation command at all, so no bibliography entry of the article is transcribed and no reference is redistributed. Reference adaptation: none was needed and none was made; every reference inside the delivered unit resolves inside it, so no unresolved reference or citation is produced. Printed slips kept literal: no printed word, symbol, hypothesis or number of the counted statement or of its proof was altered. Layout and class: the article's own class is \texttt{article} with packages including amsmath, amsthm, amssymb, amsfonts, enumerate, xcolor, authblk, appendix, hyperref, geometry, graphicx, subfigure, booktabs, multirow; the wrapper is the lane's pinned \texttt{amsart} editorial wrapper at 10pt on A4, which restates the theorem-like environments the retained text needs and adds the article's own macro definitions that the retained text needs (\texttt{\textbackslash Im}, \texttt{\textbackslash Re}), with extra wrapper packages (bm, bbm, dsfont, xspace, cancel, mathtools, enumitem). The delivered numbering is the unit's own. Assets: no retained span contains a figure environment, a graphics-inclusion command, a PDF-page inclusion, a table or a TikZ picture; no image file is copied into this package, the package declares no asset dependency and no image file is redistributed with it. Licence: version arXiv:2506.03648v1 of this article is distributed under the Creative Commons Attribution 4.0 International licence, as recorded in the arXiv licence metadata for this exact version and shown on the abstract page https://arxiv.org/abs/2506.03648v1. A full rights notice is delivered at licenses/arxiv-2506.03648-CC-BY-4.0.txt. Characters: every retained excerpt is pure ASCII, so the delivered engine, XeLaTeX with its default Latin Modern text font, reports no missing character.
\begin{equation}\label{eq-classifed-by-stokes}
\begin{cases}
\text{Type A (oscillatory solutions)}: & \Im s_{0}=1+s_{2}s_{3}>0;\\
\text{Type B (separatrix solutions)}: & \Im s_{0}=1+s_{2}s_{3}=0;\\
\text{Type C (singular solutions)}:  & \Im s_{0}=1+s_{2}s_{3}<0.
\end{cases}
\end{equation}

\begin{equation}\label{eq-stokes-full-asymptotic-case-II}
\begin{aligned}
s_{0} &\sim 2i \exp\left\{ -\sum_{s=0}^{\infty} \frac{2\Re \alpha_{s}}{\xi^{s-1}} \right\}
\cos\left( \sum_{s=0}^{\infty} \frac{2\Im \alpha_{s}}{\xi^{s-1}} \right), \\
s_{1} &= -\overline{s_{-1}} \sim i \exp\left\{ \sum_{s=0}^{\infty} \frac{2\alpha_{s}}{\xi^{s-1}} \right\}, \\
s_{2} &= -\overline{s_{3}} \sim -i \exp\left\{ \sum_{s=0}^{\infty} \frac{-4i\Im \alpha_{s}}{\xi^{s-1}} \right\}.
\end{aligned}
\end{equation}

\begin{corollary}\label{cor-classification-1}
For any fixed \( A \in (-3/2^{2/3}, C_{0}) \), there exists an increasing sequence \( \{\xi_{n}\}_{n=0}^\infty \) such that:
\begin{enumerate}[(i)]
\item The PI solutions satisfying \( y(r_n) = 0 \), \( y'(r_n) = b_n \), where \( r_n = 2A\xi_n^{4/5} \) and \( b_n = \pm 2\xi_n^{3/5} \), are of type (B);
\item For \( \xi_{2m-1} < \xi < \xi_{2m} \), \( m = 1, 2, \ldots \), the PI solutions with \( y(r) = 0 \), \( y'(r) = b \) are of type (A);
\item For \( \xi_{2m} < \xi < \xi_{2m+1} \), \( m = 0, 1, 2, \ldots \), the PI solutions with \( y(r) = 0 \), \( y'(r) = b \) are of type (C).
\end{enumerate}

Moreover, there exists \( n_{0} \in \mathbb{Z} \) such that, as \( n \to \infty \), the sequence \( \xi_{n} \) satisfies the asymptotic expansion:
\begin{equation}\label{eq-asym-xi-n}
\xi_{n} \sim \frac{\left(n + \frac{1}{2} + n_{0}\right)\pi + 2\Im \alpha_{1}}{-2\Im \alpha_{0}} \left[ 1 + \sum_{s=2}^{\infty} \frac{\mu_{s}}{\left(n - \frac{1}{2} + n_{0} + \frac{2\Im \alpha_{1}}{\pi} \right)^{s}} \right],
\end{equation}
where the coefficients \( \mu_{s} \) can be expressed in terms of \( \Im \alpha_{j} \) for \( j = 0, 1, \ldots, s \). In particular, \( \mu_{1} = -\dfrac{4\Im \alpha_{0} \Im \alpha_{2}}{\pi^{2}} \).
\end{corollary}

\begin{proof}
Fix any \( A \in (-3/2^{2/3}, C_{0}) \). From the asymptotic expression for \( s_{0} \) in Eq.~\eqref{eq-stokes-full-asymptotic-case-II}, there exists a constant \( M > 0 \) such that for all \( \xi > M \), the function \( s_{0}(\xi) \) admits a sequence of simple zeros \( \{ \xi_n \} \), viewed as a function of \( \xi \). By choosing \( \xi_1 \) appropriately, we can also arrange that \( \Im s_{0}(\xi) > 0 \) for \( \xi_{2m-1} < \xi < \xi_{2m} \) and \( \Im s_{0}(\xi) < 0 \) for \( \xi_{2m} < \xi < \xi_{2m+1} \). Comparing this behavior with the classification criterion in Eq.~\eqref{eq-classifed-by-stokes}, we establish the assertions of the corollary, except for the asymptotic formula~\eqref{eq-asym-xi-n}.


To derive the asymptotics of \( \xi_n \), we return to Eq.~\eqref{eq-stokes-full-asymptotic-case-II}. Since \( \Im \alpha_{0} < 0 \), there exists \( n_{0} \in \mathbb{Z} \) such that
\begin{equation}\label{eq-s0-cos-compare}
-2\Im \alpha_{0}\cdot\xi_{n}-2\Im \alpha_{1}-\frac{2\Im \alpha_{2}}{\xi_{n}}-\cdots\sim \left(n+n_{0}+\frac{1}{2}\right)\pi
\end{equation}
as \( n \to \infty \). Solving this asymptotic relation for \( \xi_n \), we obtain
the desired asymptotic approximation of $\xi_n$ in~\eqref{eq-asym-xi-n},
where the coefficients \( \mu_s \) can be expressed in terms of \( \Im \alpha_{j} \), for \( j = 0, 1, \ldots, s \). In particular, by substituting the expansion~\eqref{eq-asym-xi-n} into~\eqref{eq-s0-cos-compare}, we find
\[
\mu_{2}=-\frac{4\Im\alpha_{0}\Im\alpha_{2}}{\pi^2}. \qedhere
\]
\end{proof}

\end{document}
